% TODO checklist:
% - [x] Inspected src/security/audit.py (Audit event model and helpers)
% - [x] Inspected src/provenance.py (Tamper-evident provenance chain)
% - [x] Reviewed Q&A.md (Audit logging patterns)
% - [x] Reviewed README.md (Compliance features)
% - [x] Reviewed security modules for audit integration points

\section{Audit Logging}
\label{sec:audit-logging}

The Cineca Agentic Platform implements a comprehensive audit logging system that captures security-relevant events throughout the request lifecycle. The audit subsystem is designed to support compliance requirements, security forensics, and operational monitoring while maintaining performance and reliability.

\subsection{Audit Event Model}
\label{subsec:audit-event-model}

All audit events are captured using a structured data model that ensures consistency across event types:

\begin{lstlisting}[language=Python,caption={Audit event data model}]
@dataclass(slots=True)
class AuditEvent:
    event_id: str           # Unique UUID for each event
    ts: datetime            # UTC timestamp
    category: str           # auth, access, policy, ratelimit, model, data
    action: str             # login, check, allow, deny, complete, etc.
    outcome: str            # success, failure, allow, deny, info
    severity: str           # info, warning, critical
    principal: str | None   # User or service identifier
    tenant_id: str | None   # Tenant context
    resource: str | None    # Affected resource
    trace_id: str | None    # Distributed trace correlation
    meta: dict[str, Any]    # Additional context (scrubbed)
    input_hash: str | None  # SHA-256 of input (not raw content)
    output_hash: str | None # SHA-256 of output (not raw content)
\end{lstlisting}

\paragraph{Key Design Decisions:}
\begin{itemize}
    \item \textbf{Content Hashing}: Input and output are stored as cryptographic hashes, not raw content. This enables verification and correlation while preventing sensitive data from appearing in logs.
    \item \textbf{Metadata Scrubbing}: The \texttt{meta} field is automatically scrubbed of sensitive keys (passwords, tokens, secrets) before storage.
    \item \textbf{Trace Correlation}: The \texttt{trace\_id} field enables correlation with distributed tracing systems (OpenTelemetry).
\end{itemize}

\subsection{Audit Event Categories}
\label{subsec:audit-categories}

The platform categorizes audit events into distinct domains:

\begin{table}[ht]
\centering
\caption{Audit Event Categories}
\label{tab:audit-categories}
\begin{tabular}{llp{6cm}}
\hline
\textbf{Category} & \textbf{Actions} & \textbf{Description} \\
\hline
\texttt{auth} & login, logout & Authentication attempts and session lifecycle \\
\texttt{access} & allow, deny & Resource access decisions \\
\texttt{policy} & check, evaluate & Policy evaluation results \\
\texttt{ratelimit} & check & Rate limit enforcement decisions \\
\texttt{model} & complete & LLM invocations with token usage \\
\texttt{data} & read, write, delete, export & Data access operations \\
\hline
\end{tabular}
\end{table}

\subsection{Convenience Helpers}
\label{subsec:audit-helpers}

The audit module provides specialized helpers for common event types:

\subsubsection{Authentication Events}

\begin{lstlisting}[language=Python,caption={Authentication audit helpers}]
def audit_auth_success(
    *,
    username: str,
    scopes: Iterable[str] | None = None,
    tenant_id: str | None = None,
    trace_id: str | None = None,
) -> AuditEvent:
    return audit_event(
        category="auth",
        action="login",
        outcome="success",
        severity="info",
        principal=username,
        tenant_id=tenant_id,
        trace_id=trace_id,
        meta={"scopes": list(scopes or [])}
    )

def audit_auth_failure(
    *, 
    username: str | None, 
    reason: str,
    tenant_id: str | None = None
) -> AuditEvent:
    return audit_event(
        category="auth",
        action="login",
        outcome="failure",
        severity="warning",
        principal=username,
        tenant_id=tenant_id,
        meta={"reason": reason}
    )
\end{lstlisting}

\subsubsection{Access Control Events}

\begin{lstlisting}[language=Python,caption={Access control audit helper}]
def audit_access(
    *,
    principal: str | None,
    resource: str,
    method: str = "GET",
    allowed: bool,
    reason: str | None = None,
    tenant_id: str | None = None,
) -> AuditEvent:
    return audit_event(
        category="access",
        action="allow" if allowed else "deny",
        outcome="allow" if allowed else "deny",
        severity="info" if allowed else "warning",
        principal=principal,
        resource=resource,
        tenant_id=tenant_id,
        meta={"method": method, "reason": reason}
    )
\end{lstlisting}

\subsubsection{Policy Decision Events}

\begin{lstlisting}[language=Python,caption={Policy decision audit helper}]
def audit_policy_decision(
    *,
    policy: str,
    subject: str | None,
    action: str,
    resource: str,
    allowed: bool,
    reason: str | None = None,
    attributes: dict[str, Any] | None = None,
) -> AuditEvent:
    return audit_event(
        category="policy",
        action=action,
        outcome="allow" if allowed else "deny",
        severity="info" if allowed else "warning",
        principal=subject,
        resource=resource,
        meta={
            "policy": policy, 
            "reason": reason, 
            "attributes": _scrub_meta(attributes or {})
        }
    )
\end{lstlisting}

\subsubsection{Rate Limit Events}

\begin{lstlisting}[language=Python,caption={Rate limit audit helper}]
def audit_rate_limit(
    *,
    principal: str | None,
    key: str,
    allowed: bool,
    limit: int,
    window_seconds: int,
    count: int,
) -> AuditEvent:
    return audit_event(
        category="ratelimit",
        action="check",
        outcome="allow" if allowed else "deny",
        severity="info" if allowed else "warning",
        principal=principal,
        resource=key,
        meta={
            "limit": limit, 
            "window_seconds": window_seconds, 
            "count": count
        }
    )
\end{lstlisting}

\subsubsection{Model Usage Events}

\begin{lstlisting}[language=Python,caption={Model usage audit helper}]
def audit_model_usage(
    *,
    principal: str | None,
    model: str,
    prompt_tokens: int,
    completion_tokens: int,
    total_tokens: int,
    success: bool = True,
) -> AuditEvent:
    return audit_event(
        category="model",
        action="complete",
        outcome="success" if success else "failure",
        principal=principal,
        resource=model,
        meta={
            "usage": {
                "prompt_tokens": prompt_tokens,
                "completion_tokens": completion_tokens,
                "total_tokens": total_tokens
            }
        }
    )
\end{lstlisting}

\subsubsection{Data Access Events}

\begin{lstlisting}[language=Python,caption={Data access audit helper}]
def audit_data_access(
    *,
    principal: str | None,
    operation: str,  # read, write, delete, export
    data_classification: str = "internal",
    resource: str | None = None,
    record_count: int | None = None,
) -> AuditEvent:
    return audit_event(
        category="data",
        action=operation,
        outcome="success",
        principal=principal,
        resource=resource,
        meta={
            "classification": data_classification, 
            "record_count": record_count
        }
    )
\end{lstlisting}

\subsection{Sensitive Data Handling}
\label{subsec:audit-data-handling}

The audit system implements strict controls on sensitive data:

\subsubsection{Metadata Scrubbing}

\begin{lstlisting}[language=Python,caption={Sensitive metadata scrubbing}]
_SENSITIVE_KEYS = {
    "password", "authorization", "authorization_header",
    "auth", "token", "id_token", "access_token", 
    "refresh_token", "api_key", "apikey", 
    "secret", "client_secret", "ssn"
}

def _scrub_meta(meta: dict[str, Any] | None) -> dict[str, Any]:
    """Redact sensitive keys and trim large values."""
    if not meta:
        return {}
    out = {}
    for k, v in meta.items():
        lk = str(k).lower()
        if lk in _SENSITIVE_KEYS:
            out[k] = "***"
            continue
        if isinstance(v, str) and len(v) > 512:
            out[k] = v[:512] + "..."
        else:
            out[k] = v
    return out
\end{lstlisting}

\subsubsection{Content Hashing}

\begin{lstlisting}[language=Python,caption={Content hashing for audit records}]
def _sha256_hex(data: Any) -> str:
    """Generate SHA-256 hash of content for audit purposes."""
    if isinstance(data, (bytes, bytearray)):
        b = bytes(data)
    elif isinstance(data, str):
        b = data.encode("utf-8")
    else:
        b = _canonical(data).encode("utf-8")
    return hashlib.sha256(b).hexdigest()
\end{lstlisting}

\subsection{Output Channels}
\label{subsec:audit-channels}

Audit events are emitted through multiple channels for redundancy and integration:

\subsubsection{Structured Logging}

All audit events are written as structured log entries:

\begin{lstlisting}[language=Python,caption={Structured log emission}]
try:
    logger.info(
        "security_audit",
        **ev.to_dict()
    )
except Exception:
    pass  # Never fail the request path due to logging
\end{lstlisting}

The structured format enables:
\begin{itemize}
    \item Integration with log aggregation systems (ELK, Splunk, Loki)
    \item Automated alerting based on event patterns
    \item Long-term storage for compliance requirements
\end{itemize}

\subsubsection{Prometheus Metrics}

Audit events increment Prometheus counters:

\begin{lstlisting}[language=Python,caption={Prometheus metrics for audit events}]
AUDIT_EVENTS = Counter(
    "security_audit_events_total",
    "Number of security audit events",
    labelnames=("category", "action", "outcome", "severity")
)

if AUDIT_EVENTS is not None:
    AUDIT_EVENTS.labels(
        category=ev.category, 
        action=ev.action, 
        outcome=ev.outcome, 
        severity=ev.severity
    ).inc()
\end{lstlisting}

This enables:
\begin{itemize}
    \item Real-time dashboards showing security event rates
    \item Alerting on unusual patterns (e.g., spike in auth failures)
    \item SLA monitoring for security operations
\end{itemize}

\subsubsection{Provenance Chain}

Audit events are mirrored to a tamper-evident provenance chain:

\begin{lstlisting}[language=Python,caption={Provenance chain integration}]
if record_provenance is not None:
    with suppress(Exception):
        record_provenance(
            actor="security",
            action=f"audit.{ev.category}.{ev.action}",
            resource=ev.resource or "security",
            input={
                "principal": ev.principal, 
                "tenant_id": ev.tenant_id, 
                "meta": ev.meta
            },
            output={
                "outcome": ev.outcome, 
                "severity": ev.severity
            },
            meta={"trace_id": ev.trace_id} if ev.trace_id else {}
        )
\end{lstlisting}

The provenance chain provides:
\begin{itemize}
    \item Cryptographic proof of event sequence
    \item Tamper detection for audit records
    \item Long-term archival with integrity verification
\end{itemize}

\subsection{Reliability Guarantees}
\label{subsec:audit-reliability}

The audit system is designed for high reliability:

\paragraph{Fail-Safe Design:} Audit functions never raise exceptions that would interrupt the request flow. All emissions are wrapped in exception handlers:

\begin{lstlisting}[language=Python]
try:
    logger.info("security_audit", **ev.to_dict())
except Exception:
    pass  # Audit failure must not break requests
\end{lstlisting}

\paragraph{Append-Only Storage:} Audit records are written to append-only storage (log files, database tables with INSERT-only access) to prevent tampering.

\paragraph{Best-Effort Delivery:} While the system makes every effort to record all events, it prioritizes availability over audit completeness in extreme failure scenarios.

\subsection{Compliance Support}
\label{subsec:audit-compliance}

The audit system supports common compliance requirements:

\begin{itemize}
    \item \textbf{GDPR}: Data access logging with principal identification enables fulfilling access request audits.
    \item \textbf{SOC 2}: Comprehensive authentication and authorization logging supports control objectives.
    \item \textbf{HIPAA}: Audit trails for data access support breach investigation requirements.
    \item \textbf{PCI-DSS}: Logging of all access to cardholder data environments (when applicable).
\end{itemize}

\subsection{Retention and Archival}
\label{subsec:audit-retention}

Audit records are subject to configurable retention policies:

\begin{itemize}
    \item \textbf{Hot Storage}: Recent events (7-30 days) in high-performance storage for real-time querying.
    \item \textbf{Warm Storage}: Historical events (30-365 days) in cost-optimized storage with slower query performance.
    \item \textbf{Cold Storage}: Long-term archives (1+ years) in compliance with regulatory requirements.
\end{itemize}

The provenance chain provides cryptographic verification that archived records have not been modified since creation.

